The original recovery evidence did not establish the specified interval
requirements for the 50-cell reporting rule. Direct evaluation reveals
cohort and aggregation sensitivity and severe low-count null rejection, while the biological controls fail to
separate the intended mechanisms. These findings concern a documented,
limited application and its draw-dependent benchmark. They identify why its
validation claim should be withdrawn, rather than a replacement clinical
threshold.

The estimator identities explain what the original curve measures; they add
no new statistical guarantee. The practical consequence is to specify the
target and sampling unit, evaluate null rejection and effect detection with
the same interval, and report abstention and conditional performance. Any
revised rule then needs evaluation on independent draws and, for patient-level
claims, independent patients.
